These checks were decided after the main results were known and are post hoc (stages 29 and 30 of the log).

\textit{Larger faces by post count.} Scored only on face patches of at least 300 posts, the slip model alone succeeds on {{VAL:ps_E1_I30_none_300}} of the 30\,m routes and interpolation on {{VAL:ps_E1_I30_BICUBICcal_300}}; on 10\,m maps the values are {{VAL:ps_E1_I10_none_300}} and {{VAL:ps_E1_I10_BICUBICcal_300}}. Without Dumont, whose window holds rocky hills, the 30\,m values are {{VAL:ps_E1_I30_none_300_nd}} and {{VAL:ps_E1_I30_BICUBICcal_300_nd}}.

\textit{Other face rules.} A hysteresis rule (posts of 25\,degrees or more, groups of at least 10 posts holding one post of 30\,degrees or more) gives {{VAL:lab30_H_n}} patches against {{VAL:lab30_registered_n}} with the registered rule, and a rule that closes gaps of one post before grouping gives {{VAL:lab30_M_n}}. Mission success on I30 (E1) scored against these rules: slip model alone {{VAL:vb_E1_I30_none_successH}} and {{VAL:vb_E1_I30_none_successM}}, interpolation {{VAL:vb_E1_I30_BICUBICcal_successH}} and {{VAL:vb_E1_I30_BICUBICcal_successM}}.

\textit{Exposure.} Metres on face patches per km of route on I30 (E1): slip model alone {{VAL:vb_E1_I30_none_facekm}}, interpolation {{VAL:vb_E1_I30_BICUBICcal_facekm}}, boosting with Sentinel-2 {{VAL:vb_E1_I30_DEMplusS2_facekm}}, face map oracle {{VAL:vb_E1_I30_FACEMAP_facekm}}; on I10 {{VAL:vb_E1_I10_none_facekm}}, {{VAL:vb_E1_I10_BICUBICcal_facekm}} and {{VAL:vb_E1_I10_FACEMAP_facekm}}.

\textit{Survey face maps.} The face maps from simulated surveys fail because the 28\,degree rule misses faces on a coarser grid: the 2\,m survey marks {{VAL:cov2_min}} to {{VAL:cov2_max}} of the face patch posts, depending on the site, and the 5\,m survey {{VAL:cov5_min}} to {{VAL:cov5_max}}, below 60\,\% at {{VAL:cov5_below60}} sites. A survey face map needs a threshold calibrated for its own spacing.

\textit{GLO-30 alignment and migration.} At Dumont and Kelso, where GLO-30 matches the lidar without shift, interpolation of GLO-30 on its own grid reaches an AUC of {{VAL:mig_det_BICUBIC_still}}, against {{VAL:mig_det_BICUBIC_other}} at the other four sites. Applying the best shift at every site changes the mean AUC of interpolation only from {{VAL:gs_G30_BICUBIC}} to {{VAL:gs_G30shifted_BICUBIC}}, and planning with it from {{VAL:gsp_G30_BICUBICcal}} to {{VAL:gsp_G30shifted_BICUBICcal}} (slip model alone {{VAL:gsp_G30_none}} and {{VAL:gsp_G30shifted_none}}). The contrast between the two groups of sites therefore probably reflects how hard the sites are rather than misregistration.

\textit{Extended weight grid.} With the weight policy over $w$ up to 1000, held out success on I10 (E1) is {{VAL:xw_E1_I10_BICUBICcal}} for interpolation (registered grid {{VAL:xw_E1_I10_BICUBICcal_100}}), {{VAL:xw_E1_I10_GBM}} for boosting ({{VAL:xw_E1_I10_GBM_100}}), and on I30 {{VAL:xw_E1_I30_FACEMAP}} for the oracle ({{VAL:xw_E1_I30_FACEMAP_100}}); no computable layer on I30 changes by more than two points.

\textit{Finer and learned inputs.} On 5\,m block means with weights up to 1000, interpolation reaches {{VAL:pb_I5_BICUBICcal}} against {{VAL:pb_I5_none}} for the slip model alone and the oracle {{VAL:pb_I5_FACEMAP}}. On 30\,m maps the best layers from elevation alone add little: super resolution {{VAL:pb_I30_SRcal_d}} points (interval {{VAL:pb_I30_SRcal_lo}} to {{VAL:pb_I30_SRcal_hi}}) and the network trained for 6000 steps {{VAL:pb_I30_CNNlong_d}}.

\textit{Displaced face maps.} With the weight chosen on the other five sites from a grid up to 1000, where the oracle itself reaches {{VAL:pc_I30_FACEMAP}}, the face map gives {{VAL:pc_I30_shift5}} displaced by 5\,m, {{VAL:pc_I30_shift10}} by 10\,m and {{VAL:pc_I30_shift30}} by 30\,m.
